\chapter{函数的比较}\label{ch:b2:comparison}

渐近分析 --- 以简单量加可控误差代替复杂量的艺术，
已在本科一年级分册中借助泰勒展开起步。
本章将其发展为一门独立的学科：
沿一般尺度的展开、级数与积分的比较，
其完整的渐近力量，斯特林公式（完整证明），
以及对隐式定义数列的系统研究。
这些技术是渐近分析的日常工具，
此后凡估计任何量的章节 ---
级数、积分、概率 --- 都要从这张表中取用。

\section{比较关系与尺度}

\begin{definition}\label{def:b2:comparison:landau}
在点 $a$ 附近（$a \in \R$ 或 $\pm\infty$），对函数（或
数列，$n \to \infty$）：$f = o(g)$、$f = O(g)$、$f \sim g$
与本科一年级分册相同。称\emph{比较尺度}\index{比较尺度}
为 $a$ 处的一族正函数，
两两可比，且按 $o(\cdot)$ 全序排列 ---
在 $+\infty$ 处的标准尺度为
\[
x^{\alpha} (\ln x)^{\beta}
\qquad (\alpha, \beta \in \R),
\]
按 $(\alpha, \beta)$ 作字典序排列，必要时
以指数函数 $\eu^{\gamma x}$ 加细。
\end{definition}

\begin{definition}[渐近展开]\label{def:b2:comparison:expansion}
$f$ 具有如下
\emph{渐近展开}\index{渐近展开}：
\[
f = c_1 \varphi_1 + c_2\varphi_2 + \dots + c_k \varphi_k +
o(\varphi_k)
\qquad (\varphi_{i+1} = o(\varphi_i) \text{ 在该尺度中})
\]
当逐次余项满足上述估计时。
于是系数唯一：$c_1 = \lim f/\varphi_1$，
并归纳地，$c_{i+1} = \lim\,(f - \sum_{j \leq i}
c_j\varphi_j)/\varphi_{i+1}$。
\end{definition}

\begin{example}
泰勒展开是沿尺度 $(x -
a)^k$（在 $a$ 处）的渐近展开。但此概念严格更广。在 $+\infty$ 处，
\[
\frac{1}{x - \ln x}
= \frac1x \cdot \frac{1}{1 - \frac{\ln x}{x}}
= \frac1x + \frac{\ln x}{x^2} + o\Bigl(\frac{\ln x}{x^2}\Bigr),
\]
这是沿混合尺度的展开 --- 此处不适用泰勒定理，
只能用几何级数展开与 $o$ 的运算。
\end{example}

\begin{example}[标准尺度确实有序]\label{ex:b2:comparison:scaleorder}
\cref{def:b2:comparison:landau}中关于字典序的断言
需对每种情形各用一行证明。比较
$x^{\alpha}(\ln x)^{\beta}$ 与 $x^{\alpha'}(\ln x)^{\beta'}$
在 $+\infty$ 处。若 $\alpha < \alpha'$，则比值为
$x^{\alpha - \alpha'}(\ln x)^{\beta - \beta'} \to 0$，因为
$x$ 的负幂次压制任意幂的 $\ln x$（令 $x =
\eu^t$：$\eu^{(\alpha - \alpha')t}\,t^{\beta - \beta'} \to 0$
由本科一年级分册中“指数压倒多项式”的极限）。
若 $\alpha = \alpha'$ 且 $\beta < \beta'$，
则比值为 $(\ln x)^{\beta - \beta'} \to 0$，直接可得。于是数对
$(\alpha, \beta)$ 按字典序排列，
便按 $o(\cdot)$ 对尺度排序 --- 而代换 $x = \eu^t$ 是
处理混合幂与对数比较的通用技巧。
\end{example}

\begin{example}[为一群“怪兽”排序]\label{ex:b2:comparison:ranking}
尺度必须\emph{有序}；下面是一道标准的练习。
在 $+\infty$ 处比较 $n^{10}$、$\eu^{\sqrt{\ln n}\,\cdot\,\sqrt n}$、
$2^n$ 与 $n^{\ln n}$，取对数即可：
\[
10\ln n
\;\ll\; (\ln n)^2
\;\ll\; \sqrt{n\ln n}
\;\ll\; n\ln 2 ,
\]
其中 $a_n \ll b_n$ 表示 $a_n = o(b_n)$；第二项是
$\ln(n^{\ln n})$。指数运算保持这些严格的差距（若
$\ln u_n - \ln v_n \to -\infty$，则 $u_n/v_n \to 0$），故
\[
n^{10} = o\bigl(n^{\ln n}\bigr),
\qquad
n^{\ln n} = o\bigl(\eu^{\sqrt{n\ln n}}\bigr),
\qquad
\eu^{\sqrt{n\ln n}} = o(2^n) .
\]
寓意有两层：总是借助对数来比较
（比较对数的差，而非对数的比）。并且绝不要断言 $u_n \sim v_n$，
仅因 $\ln u_n \sim \ln v_n$ 成立 --- 数对 $n^{10}$
与 $n^{\ln n}$ 的 $\ln$ 比趋于 $\infty$，而 $2^n$
与 $4^n$ 的 $\ln$ 比恰为 $2$，二者却
彼此极不等价。
\end{example}

\section{级数与积分的渐近比较}

\begin{theorem}\label{thm:b2:comparison:seriesintegral}
设 $f$ 在
$\intco{1}{+\infty}$ 上连续、为正且\emph{递减}。
\begin{enumerate}
  \item 若 $\int_1^{\infty} f$ 收敛，则余项满足
        \[
        \int_{n+1}^{\infty} f \;\leq\; \sum_{k > n} f(k) \;\leq\;
        \int_{n}^{\infty} f .
        \]
  \item 若 $\int_1^\infty f$ 发散，则部分和满足
        $\sum_{k=1}^{n} f(k) = \int_1^n f + C + o(1)$，其中
        $C$ 为某常数：差 $\sum_{k \leq n} f(k) -
        \int_1^n f$ \emph{收敛}。
\end{enumerate}
\end{theorem}

\begin{proof}
由 $f(k+1) \leq \int_k^{k+1} f \leq f(k)$（递减）所作的
夹逼，已是本科一年级分册中的手法；对 $k \geq n+1$、$k \geq n$ 分别求和，
可得（1）。对（2），令 $u_k = f(k) - \int_k^{k+1} f$：
由夹逼，$0 \leq u_k \leq f(k) - f(k+1)$，故 $\sum u_k$ 的部分和
有界，其界为裂项和 $f(1) - f(n+1) \leq
f(1)$，故级数收敛。又数列 $\bigl(\int_n^{n+1}
f\bigr)_n$ 单调不增（$f$ 递减）且非负，
故收敛。记
\[
\sum_{k=1}^{n} f(k) - \int_1^n f
= \sum_{k=1}^{n} u_k + \int_n^{n+1} f ,
\]
当 $n \to \infty$ 时右端收敛，因此差
收敛于常数 $C$，即结论（2）。
\end{proof}

\begin{example}[调和数的展开]\label{ex:b2:comparison:harmonic}
设 $f(t) = \frac1t$，则 $H_n = \ln n + \gamma + o(1)$，
重新得到欧拉常数（本科一年级分册），证明更简洁。再多展开
一阶（见\cref{exo:b2:comparison:3}）：
\[
H_n = \ln n + \gamma + \frac{1}{2n} + o\Bigl(\frac1n\Bigr).
\]
数值上可清楚看出改进。取 $n = 10$ 时，$H_{10} =
2.928968\dots$，$\ln 10 = 2.302585\dots$，于是
$\gamma$ 的粗略估计为 $H_{10} - \ln 10 = 0.626383$，误差为 $0.049$；
减去修正项 $\frac1{20}$ 得 $0.576383$，
与 $\gamma = 0.577216$ 仅差 $8.3\cdot10^{-4}$ ---
这正是展开式的下一项 $\frac{1}{12\cdot100}$，
正如周末问题（第 8 题）所证。
\end{example}

\begin{example}[不用斯特林公式的粗略 $\ln(n!)$]\label{ex:b2:comparison:lnfactorial}
仅凭夹逼手法便足以估出 $\ln(n!)$。由于
$\ln$ 递增，故
\[
\int_{k-1}^{k}\ln t\,\dd t \;\leq\; \ln k \;\leq\;
\int_{k}^{k+1}\ln t\,\dd t ,
\]
对 $k = 2, \dots, n$ 求和（其中 $\int_1^n\ln = n\ln n
- n + 1$）：
\[
n\ln n - n + 1 \;\leq\; \ln(n!) \;\leq\; (n+1)\ln(n+1) - n .
\]
两界均为 $n\ln n - n + O(\ln n)$，从而 $\ln(n!) = n\ln
n - n + O(\ln n)$，特别地 $\ln(n!) \sim n\ln n$。
斯特林公式补充的是接下来的两项 --- $\frac12\ln n$
与常数 $\ln\sqrt{2\pi}$ ---，二者都依赖
\cref{thm:b2:comparison:stirling}中更精细的裂项。知道
每种工具能换得何种精度，本身就是渐近技巧的一半。
\end{example}

\begin{example}[相消须借助展开]\label{ex:b2:comparison:cancellation}
求 $\sqrt{n^2 + n} - n$ 的极限。两项均
$\sim n$，而“$\sim n - n$”毫无意义：等价量
不能相减。改用展开：
\[
\sqrt{n^2 + n} - n
= n\Bigl(\sqrt{1 + \tfrac1n} - 1\Bigr)
= n\Bigl(\frac{1}{2n} - \frac{1}{8n^2} +
O\Bigl(\frac1{n^3}\Bigr)\Bigr)
= \frac12 - \frac{1}{8n} + O\Bigl(\frac1{n^2}\Bigr) :
\]
极限为 $\frac12$，逼近速度为 $\frac1{8n}$，作为
附带结果。这一机制值得命名：两个大的
等价量之差完全由其\emph{下一项}决定，
因此须展开到两侧首次不同的那一阶 ---
并保留余项，以确认在该阶上
别无其他项残留。
\end{example}

\begin{example}[发散比较的具体计算]\label{ex:b2:comparison:loglog}
设 $f(t) = \frac{1}{t\ln t}$，定义在 $\intco{2}{+\infty}$ 上
（连续、正、递减）：$\int_2^x f = \ln\ln x -
\ln\ln 2 \to \infty$，故由
\cref{thm:b2:comparison:seriesintegral}~（2）可得
\[
\sum_{k=2}^{n}\frac{1}{k\ln k} = \ln\ln n + C + o(1)
\]
其中 $C$ 为某常数。两点启示：第一，发散
是真实的，但极其缓慢：部分和首次超过 $4$ 约在 $n
\approx \eu^{\eu^{4 - C}}$ 时，大得惊人。第二，
\emph{形状} $\ln\ln n$ 由一个原函数给出，而非
猜测得出的：对单调项，积分是标准的
求和工具，而常数 $C$ --- 与欧拉常数 $\gamma$ 一样，
--- 是对开头各项的记忆。
\end{example}

\section{斯特林公式}

\begin{lemma}[沃利斯积分，再探]\label{lem:b2:comparison:wallis}
设 $W_n = \int_0^{\pi/2} \sin^n t\,\dd t$，则 $nW_nW_{n-1} =
\frac\pi2$（$n \geq 1$），$(W_n)$ 递减，且 $W_n \sim
\sqrt{\dfrac{\pi}{2n}}$。
\end{lemma}

\begin{proof}
分部积分得 $nW_n = (n-1)W_{n-2}$（$n \geq 2$），于是
$nW_nW_{n-1}$ 关于 $n$ 为常数，等于 $1 \cdot W_1 W_0 =
\frac\pi2$。递减性：$\sin^{n+1} \leq \sin^n$ 在
$\intcc{0}{\frac\pi2}$ 上成立。夹逼的细节：单调性
给出 $W_{n+1} \leq W_n \leq W_{n-1}$，两边除以 $W_{n-1}
> 0$，
\[
\frac{n}{n+1} = \frac{W_{n+1}}{W_{n-1}} \leq
\frac{W_n}{W_{n-1}} \leq 1 ,
\]
左端的等式由指标为 $n + 1$ 的递推式得出。两个
界都趋于 $1$：$W_n \sim W_{n-1}$，由此得
\[
nW_n^2 \sim nW_nW_{n-1} = \frac\pi2
\qquad\Longrightarrow\qquad
W_n \sim \sqrt{\frac{\pi}{2n}} .
\]
\end{proof}

\begin{example}[最初几个沃利斯积分]\label{ex:b2:comparison:wallisvalues}
由 $W_0 = \frac\pi2$、$W_1 = 1$ 及递推式 $nW_n =
(n-1)W_{n-2}$：
\[
W_2 = \frac\pi4, \qquad
W_3 = \frac23, \qquad
W_4 = \frac{3\pi}{16}, \qquad
W_5 = \frac{8}{15}, \qquad
W_6 = \frac{5\pi}{32}.
\]
偶数下标的项含 $\pi$，奇数下标的项为有理数 --- 两者
恰是闭式的两条交错乘积链。数值上 $W_6
\approx 0.4909$，对照渐近式
$\sqrt{\pi/12} \approx 0.5116$：在 $n = 6$ 时，等价关系
已在 $5\%$ 之内，且乘积恒等式对每个
$n$ 都精确成立：$6\,W_6W_5 = 6\cdot\frac{5\pi}{32}\cdot\frac8{15} =
\frac\pi2$。像这样的小表格，是抓出代数笔误最省力的办法，
趁笔误尚未污染渐近论证时即可发现。
\end{example}

\begin{theorem}[斯特林]\label{thm:b2:comparison:stirling}
\[
n! \;\sim\; \sqrt{2\pi n}\, \Bigl(\frac{n}{\eu}\Bigr)^{\!n}
.\index{Stirling's formula}
\]
\end{theorem}

\begin{proof}
\emph{第一步：$n! \sim C \sqrt n\, (n/\eu)^n$，其中常数 $C >
0$。} 设
\[
d_n = \ln(n!) - \Bigl(n + \frac12\Bigr)\ln n + n .
\]
则
\[
d_n - d_{n+1}
= \Bigl(n + \frac12\Bigr) \ln\frac{n+1}{n} - 1
= \Bigl(n + \frac12\Bigr)\Bigl(\frac1n - \frac{1}{2n^2} +
\frac{1}{3n^3} + o\bigl(n^{-3}\bigr)\Bigr) - 1
= \frac{1}{12n^2} + o\Bigl(\frac{1}{n^2}\Bigr),
\]
由 $\ln(1 + \frac1n)$ 的泰勒展开得出。级数 $\sum (d_n
- d_{n+1})$ 绝对收敛（与 $\sum
n^{-2}$ 比较），故 $(d_n)$ 收敛，记其极限为 $d$；取指数得 $n! \sim
C\sqrt n\,(n/\eu)^n$，其中 $C = \eu^{d}$。

\emph{第二步：借助沃利斯公式求 $C = \sqrt{2\pi}$。} 闭式 $W_{2p}
= \frac{(2p)!}{4^p (p!)^2}\cdot\frac\pi2$（由递推式得到，
一年级的计算在 \cref{lem:b2:comparison:wallis} 的
设定下重做），与第一步结合：
\[
W_{2p} \sim \frac{C\sqrt{2p}\,(2p/\eu)^{2p}}
{4^p\,\bigl(C\sqrt p\,(p/\eu)^p\bigr)^2}\cdot\frac{\pi}{2}
= \frac{\sqrt{2p}}{C\,p}\cdot\frac{\pi}{2}
= \frac{\pi}{C}\cdot\frac{1}{\sqrt{2p}} .
\]
比较 $W_{2p} \sim \sqrt{\frac{\pi}{4p}}$
（见\cref{lem:b2:comparison:wallis}）：$\frac{\pi}{C\sqrt{2p}} =
\sqrt{\frac{\pi}{4p}}\,(1 + o(1))$，可得 $C = \pi
\sqrt{\frac{4p}{2p\,\pi}} = \sqrt{2\pi}$。
\end{proof}

\begin{example}[中心二项式系数]\label{ex:b2:comparison:centralbinomial}
\[
\binom{2n}{n} = \frac{(2n)!}{(n!)^2}
\sim \frac{\sqrt{4\pi n}\,(2n/\eu)^{2n}}{2\pi n\,(n/\eu)^{2n}}
= \frac{4^n}{\sqrt{\pi n}} :
\]
对称随机游走回到 $0$ 的概率，
在时刻 $2n$ 为 $\sim \frac{1}{\sqrt{\pi n}}$ --- 对
\cref{ch:b2:randomvar} 的预告。
\end{example}

\begin{remark}[本卷内的展望]\label{rem:b2:comparison:perspectives}
此后每一章的定量内容都沿用本章的
语言。\cref{ch:b2:series} 通过将各项与尺度
$n^{-\alpha}(\ln n)^{-\beta}$ 比较来给级数分类 --- 其
周末问题完整刻画了这一边界。
\cref{ch:b2:integration} 对反常积分做了同样的事，
在连续变量中使用完全相同的尺度。
\cref{ch:b2:powerseries} 计算收敛半径，其依据是
$\limsup\abs{a_n}^{1/n}$，这是一道 $n$ 次方根的等价练习，
其标准关键步骤是斯特林公式（$\sqrt[n]{n!}
\sim \frac n\eu$，见\cref{exo:b2:comparison:4}）。而概率
各章直接使用斯特林公式：二项式系数的局部
估计在\cref{ch:b2:randomvar}中，即
\cref{ex:b2:comparison:centralbinomial} 与
\cref{ex:b2:comparison:lopsided} 的原样。渐近
分析不是这里的一章，而是本卷的特色。
\end{remark}

\begin{method}[逐级代入检查清单]\label{met:b2:comparison:checklist}
在相信逐级代入得到的展开之前，请检查以下四点。
(1) \emph{先证存在性：}根或数列须首先被
确定（单调性、介值性），然后才能展开
--- 无所指的符号展开得很漂亮，
毫无意义。(2) \emph{每趟只算一阶：}每次代入
只能相信到所输入估计的阶；从一趟中同时提取
两个新项，是造成错误系数的经典来源。
(3) \emph{余项随行：}每一步代数运算都要携带
$o(\cdot)$，并让吸收
（较小项被较大余项吞没）
在最后明确地完成。(4) \emph{数值检验：}在一个
可信的 $n$ 值处求值：系数错误常能经受住代数
重新推导的考验，却几乎经不住
算术的检验。
\end{method}

\begin{remark}[常见陷阱]\label{rem:b2:comparison:pitfalls}
(i) 等价式\emph{相加}时很不可靠：由 $u_n \sim n + \ln n$ 与
$v_n \sim -n$，\emph{不能}推出 $u_n + v_n \sim \ln
n$；抵消需要带显式余项的展开，
绝不能只用裸的等价式。(ii) 等价式绝不可取指数：
$n + 1 \sim n$，但 $\eu^{n+1} \not\sim \eu^n$；安全的
做法是对趋于
$+\infty$ 的等价式取对数（本章周末问题第 24 题）。(iii)
渐近展开依附于一个\emph{尺度}：写作
$f = \frac1x + o\bigl(\frac1{x^2}\bigr)$ 所断言的内容多于 $f =
\frac1x + o\bigl(\frac1x\bigr)$，混用二者会使
后续代数失效。(iv) 逐级代入时，应代入
\emph{整个}当前展开，包括余项 --- 中途丢掉
某个 $o(\cdot)$，会得到看似合理实则错误的
系数。(v) 级数与积分的比较需要单调性：对振荡项
它完全失效（比较
$\sum\frac{\sin k}k$，见\cref{ch:b2:series}）。
\end{remark}

\begin{example}[斯特林公式的数值检验]\label{ex:b2:comparison:stirlingnum}
当 $n = 10$ 时，公式给出 $\sqrt{20\pi}\,(10/\eu)^{10}
\approx 3\,598\,696$，与 $10! = 3\,628\,800$ 相比，相对误差为
$8.3\cdot10^{-3}$，就“渐近”论断而言，在
$n = 10$ 处已相当可观。误差具有一定结构 --- 精确的加细式 $n!
= \sqrt{2\pi n}\,(n/\eu)^n\bigl(1 + \frac1{12n} +
O(n^{-2})\bigr)$ --- 其首项修正 $\frac1{120} \approx
8.3\cdot10^{-3}$ 几乎恰好解释了观察到的差距。
周末问题中的欧拉–麦克劳林方法，
正是系统产生此类修正项的源泉。
\end{example}

\begin{remark}[本章的应用]
渐近比较是此后一切定量内容的“语法”，
其应用包括：
\cref{ch:b2:series}中的收敛判别法与伯特兰德全景、
\cref{ch:b2:integration}中的可积性判别准则、
\cref{ch:b2:powerseries}中的收敛半径计算，以及
\cref{ch:b2:randomvar}中的极限定理（其中斯特林公式用于
棣莫弗–拉普拉斯估计）。本科三年级分册将
本章手工证明的一条思路 --- 提取主项、控制
其余部分 --- 系统化为拉普拉斯方法与控制收敛定理。
\end{remark}

\begin{example}[与自身比较的积分：$\int_2^x \frac{\dd t}{\ln t}$]\label{ex:b2:comparison:liworked}
比较工具同样适用于积分。设 $F(x) =
\int_2^x\frac{\dd t}{\ln t}$（被积函数在
$\intco2\infty$ 上连续）。分部积分得
\[
F(x) = \Bigl[\frac{t}{\ln t}\Bigr]_2^x +
\int_2^x\frac{\dd t}{(\ln t)^2}
= \frac{x}{\ln x} + O\Bigl(\int_2^x\frac{\dd t}{(\ln
t)^2}\Bigr) + O(1),
\]
余项积分为 $o\bigl(\frac{x}{\ln x}\bigr)$：
在 $\sqrt x$ 处将其分为两段，分别放缩为
\[
\int_2^{\sqrt x}\frac{\dd t}{(\ln t)^2} \leq \sqrt x
\qquad\text{和}\qquad
\int_{\sqrt x}^{x}\frac{\dd t}{(\ln t)^2} \leq
\frac{x}{(\ln\sqrt x)^2} = \frac{4x}{(\ln x)^2} .
\]
于是 $F(x) \sim \frac{x}{\ln x}$。若读者在本章周末问题中接触过素数定理，
便会认出这一函数：
$F$ 即对数积分，是更好的
$\pi(x)$ 的估计，计算表明它
在一阶上与 $\frac{x}{\ln x}$ 相一致。
\end{example}

\begin{example}[斯特林公式用于不对称二项式系数]\label{ex:b2:comparison:lopsided}
沿用
\cref{ex:b2:comparison:centralbinomial}中的三阶乘估计步骤，对
$\binom{3n}{n} = \frac{(3n)!}{n!\,(2n)!}$ 给出：
\[
\binom{3n}{n} \sim
\frac{\sqrt{6\pi n}\,(3n/\eu)^{3n}}
{\sqrt{2\pi n}\,(n/\eu)^{n}\cdot\sqrt{4\pi n}\,(2n/\eu)^{2n}}
= \sqrt{\frac{3}{4\pi n}}\,
\Bigl(\frac{27}{4}\Bigr)^{\!n} .
\]
指数增长率 $\frac{27}4 = \frac{3^3}{2^2}$ 可记作
信息论熵记号中的 $\eu^{3n\,H(1/3)}$，
即不对称二项式的增长严格慢于中心
二项式每两步为 $4^n$ --- 此处 $(27/4)^{1/3} \approx 1.89 < 2$，即每一步。
组合学与概率论中的每一个二项式渐近式
（见\cref{ch:b2:randomvar}）都是同一计算，
只是权重不同。
\end{example}

\section{隐式定义的数列}

\begin{method}\label{met:b2:comparison:implicit}
求方程的解 $x_n$ 的渐近性态，该方程为 $F(x, n) =
0$：\index{隐式数列}
\begin{enumerate}
  \item \emph{定位}：证明 $x_n$ 在
        确定区间内存在且唯一（借助单调性与介值定理），
        并求其粗略性态（极限、
        增长阶）。
  \item \emph{迭代修正}：将粗略形式 $x_n = (\text{主
        项})(1 + \varepsilon_n)$ 代入方程，解出
        $\varepsilon_n$ 的下一阶；重复此步，每一轮
        加细一阶。
\end{enumerate}
\end{method}

\begin{example}\label{ex:b2:comparison:tan}
当 $n \geq 1$ 时，方程 $\tan x = x$ 恰有一个解，
记为 $x_n$，它落在 $\intoo{n\pi - \frac\pi2}{n\pi + \frac\pi2}$ 内（函数
$\tan x - x$ 由 $-\infty$ 增至 $+\infty$，其
导数为 $\tan^2 x \geq 0$）。\emph{粗略估计：} $x_n = n\pi +
\frac\pi2 - y_n$，其中 $y_n \in \intoo{0}{\pi}$；由于 $x_n \to
\infty$ 且 $\tan x_n = x_n \to +\infty$，故 $x_n$ 自左侧
趋近于渐近线：$y_n \to 0$。\emph{迭代修正：} $\tan x_n =
\cot y_n = \frac{1}{\tan y_n} \sim \frac{1}{y_n}$，而方程
$\cot y_n = x_n \sim n\pi$ 给出 $y_n \sim \frac{1}{n\pi}$。于是
\[
x_n = n\pi + \frac\pi2 - \frac{1}{n\pi} +
o\Bigl(\frac1n\Bigr),
\]
该过程可继续至任意阶
（见\cref{exo:b2:comparison:6}）。
\end{example}

\begin{example}[方法的第二次应用]\label{ex:b2:comparison:xplusln}
渐近地求解 $x + \ln x = n$。\emph{定位：} $x
\mapsto x + \ln x$ 由 $-\infty$ 增至 $+\infty$，在
$\intoo{0}{+\infty}$ 上恰有唯一根 $x_n$，且 $x_n \to \infty$。
\emph{粗略估计：} 由 $\ln x_n = o(x_n)$ 得 $x_n \sim n$。
\emph{迭代修正：} 由 $x_n = n - \ln x_n$ 及 $\ln x_n = \ln n
+ o(1)$（等价量取对数，两边均 $\to \infty$）：
\[
x_n = n - \ln n + o(1) ;
\]
再迭代一次，此时 $\ln x_n = \ln\bigl(n - \ln n + o(1)\bigr) =
\ln n - \frac{\ln n}{n} + o\bigl(\frac{\ln n}n\bigr)$:
\[
x_n = n - \ln n + \frac{\ln n}{n} +
o\Bigl(\frac{\ln n}{n}\Bigr).
\]
（取 $n = 100$ 检验：根约为 $x \approx 95.4415$；
三项公式给出 $100 - 4.6052 + 0.0461 = 95.4409$，
两项公式为 $95.3948$ --- 每一轮都获得预期的
阶数。）同一循环，换到第三种情形：
\cref{met:b2:comparison:implicit}并不在乎
方程的形状，只要求每一轮
都能分离出主导的未知量。
\end{example}

\section{练习}

\begin{exercise}[$\star$]\label{exo:b2:comparison:1}
在 $+\infty$ 处展开，除首项外再取两项：
\[
\sqrt{x^2 + x + 1} ,
\qquad
\ln(x^2 + x) - 2\ln x,
\qquad
\frac{x + \sin x}{x - \ln x} .
\]
\end{exercise}

\begin{exercise}[$\star$]\label{exo:b2:comparison:2}
讨论以下和式的敛散性（收敛或发散），并在发散时给出
$\sum_{k \leq n} k^\alpha$ 的首项渐近，分 $\alpha >
-1$、$\alpha = -1$、$\alpha < -1$ 三种情形讨论，借助
\cref{thm:b2:comparison:seriesintegral}。
\end{exercise}

\begin{exercise}[$\star\star$]\label{exo:b2:comparison:3}
证明 $H_n = \ln n + \gamma + \frac{1}{2n} + o\bigl(\frac1n\bigr)$。
\emph{（研究 $v_n = H_n - \ln n - \gamma$：证明 $v_n - v_{n+1} =
\frac{1}{2n^2} + O(n^{-3})$，并对尾项求和，与
$\sum_{k \geq n} \frac{1}{2k^2} \sim \frac{1}{2n}$ 作比较 ---
由\cref{thm:b2:comparison:seriesintegral}~(1)可知。）}
\end{exercise}

\begin{exercise}[$\star\star$]\label{exo:b2:comparison:4}
利用斯特林公式，求下列各式的等价量：$\dfrac{(3n)!}{(n!)^3}$；
$\;\dfrac{n!}{n^n}$；$\;\sqrt[n]{n!}$（写成 $\frac n\eu(1 + o(1))$ 的形式，
并精确到两项）。
\end{exercise}

\begin{exercise}[$\star\star$]\label{exo:b2:comparison:5}
设 $n \geq 2$，证明 $x^n + x = 1$ 恰有一个解 $x_n
\in \intoo{0}{1}$，且 $x_n \to 1$，并证明
\[
x_n = 1 - \frac{\ln n}{n} + o\Bigl(\frac{\ln n}{n}\Bigr).
\]
\emph{（由 $x_n^n = 1 - x_n$ 取对数，并利用
$x_n = 1 - \varepsilon_n$ 逐步迭代。）}
\end{exercise}

\begin{exercise}[$\star\star$]\label{exo:b2:comparison:6}
将\cref{ex:b2:comparison:tan}再推进一阶：
\[
x_n = n\pi + \frac\pi2 - \frac{1}{n\pi} + \frac{1}{2n^2\pi} +
o\Bigl(\frac{1}{n^2}\Bigr).
\]
\emph{（精确写出 $\cot y_n = x_n$，展开 $\cot y = \frac1y -
\frac y3 + o(y)$，以及 $x_n = n\pi(1 + \frac{1}{2n} - \dots)$，并
对照确定各项。）}
\end{exercise}

\begin{exercise}[$\star\star$]\label{exo:b2:comparison:7}
求 $\lim_{n\to\infty} \dfrac{1}{n!}\sum_{k=0}^{n} k!$
\emph{（对去掉最后两项后的和式作放缩）}，并推出渐近展开式
$\sum_{k \leq n} k! = n!\bigl(1 + \frac1n +
O(n^{-2})\bigr)$。
\end{exercise}

\begin{exercise}[$\star\star\star$]\label{exo:b2:comparison:8}
设 $u_0 > 0$ 且 $u_{n+1} = u_n + \dfrac{1}{u_n}$。证明 $u_n
\to \infty$，再证 $u_n \sim \sqrt{2n}$ \emph{（研究 $u_n^2$：
其增量为 $2 + u_n^{-2}$；求和）}，并进一步精确化：
\[
u_n = \sqrt{2n}\Bigl(1 + \frac{\ln n}{8n} +
o\Bigl(\frac{\ln n}{n}\Bigr)\Bigr).
\]
\emph{（由 $u_n^2 = 2n + \sum_{k<n} u_k^{-2} + u_0^2$ 及 $u_k^2
\sim 2k$，和式 $\sim \frac12\ln n$ 由
\cref{thm:b2:comparison:seriesintegral}可得。）}
\end{exercise}

\begin{exercise}[$\star\star\star$]\label{exo:b2:comparison:9}
（一个略作变形的黎曼和）求下列量的渐近性态：
\[
S_n = \sum_{k=1}^{n} \frac{1}{n + k\ln n} .
\]
\emph{（提出 $n$：$S_n = \frac1n\sum_k \bigl(1 +
\frac{k\ln n}{n}\bigr)^{-1}$；识别出一个带
缓变参数 $t = \ln n$ 的黎曼型和，计算 $\int_0^1
\frac{\dd u}{1 + tu} = \frac{\ln(1+t)}{t}$，并得出 $S_n \sim
\frac{\ln\ln n}{\ln n}$。）}
\end{exercise}

\begin{exercise}[$\star$]\label{exo:b2:comparison:10}
证明恒等式 $(\ln n)^{\ln n} = n^{\ln\ln n}$，再将
以下各式在无穷远处按 $o(\cdot)$ 由小到大排序，并给出
证明：$n^2$、$(\ln n)^{\ln n}$、$2^n$、$n!$、$n^n$。
\end{exercise}

\begin{exercise}[$\star\star$]\label{exo:b2:comparison:11}
（$\sum 1/k^2$ 的尾项，取两项）利用精确的裂项相消
$\sum_{k > n} \frac{1}{k(k+1)} = \frac{1}{n+1}$ 及
分解式 $\frac1{k^2} = \frac{1}{k(k+1)} +
\frac{1}{k^2(k+1)}$，证明
\[
\sum_{k > n} \frac{1}{k^2}
= \frac1n - \frac{1}{2n^2} + O\Bigl(\frac{1}{n^3}\Bigr).
\]
\end{exercise}

\begin{exercise}[$\star\star\star$]\label{exo:b2:comparison:12}
设 $u_0 = \frac12$ 且 $u_{n+1} = u_n + \eu^{-u_n}$。证明
$u_n \to \infty$，再 --- 令 $v_n = \eu^{u_n}$，并证明
$v_{n+1} = v_n + 1 + \frac{1}{2v_n} + O\bigl(v_n^{-2}\bigr)$ ---
由此建立如下：
\[
u_n = \ln n + \frac{\ln n}{2n} + O\Bigl(\frac1n\Bigr).
\]
\end{exercise}

\section{问题：自举法，从欧拉–麦克劳林到
质数}

隐式的或累积而成的量很少能一次给出其渐近性态；
通常需分若干轮提取，每一轮都把上一轮的估计
代回定义关系式中。
本周末问题在新方程上训练这一循环，
证明一阶的\emph{欧拉–麦克劳林公式}（即级数与积分比较的梯形改进，
附带严格的误差界），反解 $x\ln x = n$，
并兑现该方法最著名的一张“支票”：
由已承认的素数定理，
得出 $p_n \sim n\ln n$ 这一关于第 $n$ 个素数的渐近律。

\begin{problem}[{周末问题 --- 欧拉–麦克劳林修正与
第 $n$ 个素数的渐近性}]
\label{pb:b2:comparison:1}

\medskip
\noindent\textbf{第一部分 --- 在新
方程上的自举循环。}
\begin{enumerate}
  \item 证明
        \cref{def:b2:comparison:expansion}中的唯一性断言：若 $f = \sum_{i\leq
        k} c_i\varphi_i + o(\varphi_k) = \sum_{i \leq k}
        c_i'\varphi_i + o(\varphi_k)$ 沿同一尺度，则
        $c_i = c_i'$ 对所有 $i$ 成立。随后将课程中的混合
        例子再推进一级：
        \[
        \frac{1}{x - \ln x} = \frac1x + \frac{\ln x}{x^2} +
        \frac{(\ln x)^2}{x^3} + o\Bigl(\frac{(\ln
        x)^2}{x^3}\Bigr) \qquad (x \to +\infty),
        \]
        并说明为何不出现 $\frac{c}{x^2}$ 项。
  \item 证明：对每个 $n \geq 1$，方程 $\eu^x + x =
        n$ 恰有一个实根 $x_n$，并且 $x_n
        \to +\infty$，且 $x_n \sim \ln n$。
  \item 作两次自举：
        \[
        x_n = \ln n - \frac{\ln n}{n} - \frac{(\ln n)^2}{2n^2} +
        o\Bigl(\frac{(\ln n)^2}{n^2}\Bigr).
        \]
  \item 取 $n = 1000$ 进行数值检验：比较 $x_{1000}
        \approx 6.90083$ 与第 3 题中的一项、两项、三项值，
        精确到小数点后五位。
\end{enumerate}

\medskip
\noindent\textbf{第二部分 --- 一阶欧拉–麦克劳林公式。}
\begin{enumerate}[resume]
  \item 证明梯形核恒等式：设 $g$ 为
        $C^2$ 类函数，定义在 $\intcc{0}{1}$ 上，
        \[
        \int_0^1 g(t)\,\dd t = \frac{g(0) + g(1)}{2}
        - \frac12\int_0^1 t(1 - t)\,g''(t)\,\dd t
        \]
        \emph{（对 $\frac12 t(1-t)g''$ 分部积分两次）}。
  \item 设 $f$ 为 $C^2$ 类，定义在 $\intco{1}{+\infty}$ 上，且
        $\int_1^\infty \abs{f''} < \infty$。证明
        \[
        E_n = \sum_{k=1}^{n} f(k) - \int_1^n f -
        \frac{f(1) + f(n)}{2}
        \]
        收敛于常数 $E$，并满足尾项估计
        $\abs{E - E_n} \leq \frac18\int_n^\infty\abs{f''}$，
        这便是一阶的\emph{欧拉–麦克劳林公式}。
  \item 将此结果用于 $f(t) = \frac1t$，证明
        \[
        H_n = \ln n + \gamma + \frac{1}{2n} + \varepsilon_n,
        \qquad \abs{\varepsilon_n} \leq \frac{1}{8n^2},
        \]
        加强\cref{exo:b2:comparison:3}的结论（将常数与
        $\gamma$ 等同，方法是与
        \cref{ex:b2:comparison:harmonic}比较）。
  \item 求下一项系数：证明 $\varepsilon_n =
        -\frac{1}{12n^2} + o\bigl(\frac1{n^2}\bigr)$
        \emph{（$E_n$ 的增量为
        $\frac12\int_0^1t(1-t)f''(n+t)\dd t =
        \frac1{12}f''(n) + o(f''(n))$；借助
        \cref{thm:b2:comparison:seriesintegral}对尾项求和）}。
  \item 将第 6 题用于 $f = \ln$：用三行重新推导
        数列 $d_n = \ln n! - (n +
        \frac12)\ln n + n$ 的收敛性（
        \cref{thm:b2:comparison:stirling}的第 1 步），并附带
        误差阶 $d_n = d + O\bigl(\frac1n\bigr)$。
  \item 将第 6 题用于 $f(t) = \frac{1}{\sqrt t}$，证明
        \[
        \sum_{k=1}^{n}\frac1{\sqrt k} = 2\sqrt n + c +
        \frac{1}{2\sqrt n} + O\Bigl(\frac{1}{n^{3/2}}\Bigr)
        \]
        其中 $c$ 为某个常数，并在 $n =
        10^4$ 处计算各项（常数 $c \approx -1.4604$）。
\end{enumerate}

\medskip
\noindent\textbf{第三部分 --- 反解：方程 $x\ln x =
n$。}
\begin{enumerate}[resume]
  \item 证明方程 $x\ln x = n$ 恰有一个解 $x_n \in
        \intco{1}{+\infty}$，对 $n \geq 1$ 成立，且 $x_n \to
        \infty$，并且 $\ln x_n \sim \ln n$。
  \item 由此得一项反解 $x_n \sim
        \dfrac{n}{\ln n}$，再作一次自举：
        \[
        \ln x_n = \ln n - \ln\ln n + o(1),
        \qquad
        x_n = \frac{n}{\ln n}\Bigl(1 + \frac{\ln\ln n}{\ln n}
        + o\Bigl(\frac{\ln\ln n}{\ln n}\Bigr)\Bigr).
        \]
  \item 在 $n = 10^6$ 处检验：真根约为 $x \approx
        87\,848$；与一项（$\approx 72\,382$）
        及两项（$\approx 86\,140$）的近似值比较，并解释
        收敛缓慢的原因（展开参数为 $\frac{\ln\ln
        n}{\ln n}$ 仅 $\approx 0.19$，取 $n = 10^6$ 时如此）。
  \item 现在\emph{承认}素数定理：
        $\pi(x)$ 为满足 $\leq x$ 的质数个数，且 $\pi(x) \sim
        \frac{x}{\ln x}$（$x \to \infty$ 时成立，严格证明见
        本科三年级分册）。记 $p_n$ 为第 $n$ 个质数，
        说明 $\pi(p_n) = n$，并沿用
        第 11--12 题的反解方法，证明：
        \[
        p_n \sim n \ln n .
        \]
  \item 附加题：（a）证明 $\sum_{k \leq n} p_k \sim
        \frac{n^2\ln n}{2}$ \emph{（将 $\sum k\ln k$ 与
        $\int t\ln t\,\dd t$）}；（b）计算
        在 $100$ 位整数中均匀随机取一个，它为质数的近似概率
        （$\ln 10^{100} \approx 230.26$，因此约为
        $230$ 分之一）。
\end{enumerate}

\medskip
\noindent\textbf{第四部分 --- 导出的方法：$x\tan x =
1$.}
\begin{enumerate}[resume]
  \item 证明：对每个 $n \geq 1$，方程 $\tan x =
        \frac1x$ 恰有一个解 $x_n$，
        位于 $\intoo{n\pi}{\,n\pi + \frac\pi2}$ 内，且 $z_n = x_n
        - n\pi \to 0^+$.
  \item 取一项：$z_n \sim \dfrac{1}{n\pi}$。
  \item 证明 $z_n$ 的展开式\emph{不含}
        $\frac{c}{n^2}$ 项：$z_n = \frac1{n\pi} +
        O\bigl(\frac{1}{n^3}\bigr)$.
  \item 取三项：利用 $\arctan u = u - \frac{u^3}3 +
        O(u^5)$ 及 $\frac1{x_n} = \frac{1}{n\pi} -
        \frac{z_n}{(n\pi)^2} + O(n^{-3}\cdot z_n^2)$，证明
        \[
        x_n = n\pi + \frac{1}{n\pi} -
        \frac{4}{3\pi^3 n^3} + o\Bigl(\frac{1}{n^3}\Bigr).
        \]
  \item 取 $n = 3$ 检验：真根 $x_3 \approx 9.5293344$；
        比较一项与三项的近似值，并用一句话与课程中的
        $\tan x = x$ 情形对照：
        （见\cref{ex:b2:comparison:tan}），说明各数列
        位于其区间中的何处，以及原因。
\end{enumerate}

\medskip
\noindent\textbf{第五部分 --- 动力学自举、游戏
规则与综合。}
\begin{enumerate}[resume]
  \item 设 $u_0 \in \intoo{0}{\pi}$，且 $u_{n+1} = \sin u_n$。
        证明 $u_n \to 0$ 单调递减，并求
        $\dfrac{1}{u_{n+1}^2} - \dfrac{1}{u_n^2}$ 的极限\emph{（展开
        $\sin^{-2}$，借助 $\sin u = u - \frac{u^3}6 + o(u^3)$）}。
  \item 由切萨罗平均（本科一年级分册）推出经典结论：
        \[
        u_n \sim \sqrt{\frac{3}{n}} .
        \]
  \item （可验证的数值计算）利用第 7 题的严格界，
        证明：将 $\ln n + \gamma + \frac1{2n}$ 在 $n
        = 10^6$ 处求值，所得的 $H_{10^6}$ 误差至多为
        $1.25\cdot10^{-13}$ --- 百万项之和，
        仅用三项即算到十三位数字。
  \item （游戏规则）对下列陈述，证明或否定，并给出
        证明或反例：(a) 若 $u_n \sim v_n \to +\infty$，则
        $\ln u_n \sim \ln v_n$；(b) 若 $u_n \sim v_n$，则
        $\eu^{u_n} \sim \eu^{v_n}$；(c) 若 $f \sim g$ 在
        $+\infty$ 处（$f, g$ 可微），则 $f' \sim g'$。
  \item （综合）每项用一句话说明：第 I、III、IV 部分中所用的
        \cref{met:b2:comparison:implicit} 自举循环；
        梯形修正相对于
        \cref{thm:b2:comparison:seriesintegral} 增加了什么；为何
        对 $x\ln x$ 取反函数恰好是从 $\pi(x)$ 到
        $p_n$ 的桥梁；第 24 题的哪条规则保护了哪一步。
        并指出两座高峰：欧拉–麦克劳林公式（一阶）
        与第 $n$ 个
        质数的渐近律。
\end{enumerate}
\end{problem}
